\title{复数（2016--2025 高考真题）}
\author{}
\date{}
\maketitle
\noindent 本专题共 71 道题，按年份排列。每题标注原卷与原题号。

\section*{2016 年}

\begin{question}[points=5]
\textbf{来源：}2016 年 全国III卷-理科，第 2 题\par
若 $z=1+2i$, 则 $\dfrac{4i}{z\overline{z}-1}=$
\begin{choices}
  \item $1$
  \item $-1$
  \item $i$
  \item $-i$
\end{choices}
\begin{solution}
\textbf{答案：}$i$

\textbf{分析：}利用复数的乘法运算法则, 化简求解即可.

\textbf{详解：}$z=1+2i$, 则 $\dfrac{4i}{z\overline{z}-1}=\dfrac{4i}{(1+2i)(1-2i)-1}=\dfrac{4i}{1-4i^2-1}=\dfrac{4i}{4}=i$, 故选 C.
\end{solution}
\end{question}

\begin{question}[points=5]
\textbf{来源：}2016 年 全国III卷-文科，第 2 题\par
若 $z=4+3i$, 则 $\dfrac{\overline{z}}{|z|}=$ ( )
\begin{choices}
  \item $1$
  \item $-1$
  \item $\dfrac{4}{5}+\dfrac{3}{5}i$
  \item $\dfrac{4}{5}-\dfrac{3}{5}i$
\end{choices}
\begin{solution}
\textbf{答案：}$\dfrac{4}{5}-\dfrac{3}{5}i$

\textbf{分析：}利用复数的除法以及复数的模化简求解即可.

\textbf{详解：}$z=4+3i$, 则 $\dfrac{\overline{z}}{|z|}=\dfrac{4-3i}{\sqrt{4^2+3^2}}=\dfrac{4-3i}{5}=\dfrac{4}{5}-\dfrac{3}{5}i$. 故选 D.
\end{solution}
\end{question}

\begin{question}[points=5]
\textbf{来源：}2016 年 全国II卷-理科，第 1 题\par
已知$z=(m+3)+(m-1)i$在复平面内对应的点在第四象限，则实数$m$的取值范围是（\quad）
\begin{choices}
  \item $(-3,1)$
  \item $(-1,3)$
  \item $(1,+\infty)$
  \item $(-\infty,-3)$
\end{choices}
\begin{solution}
\textbf{答案：}$(-3,1)$

\textbf{分析：}利用复数对应点所在象限，列出不等式组求解即可．

\textbf{详解：}由题意得$\begin{cases} m+3>0 \\ m-1<0 \end{cases}$，解得$-3<m<1$．故选A．
\end{solution}
\end{question}

\begin{question}[points=5]
\textbf{来源：}2016 年 全国II卷-文科，第 2 题\par
设复数$z$满足$z+i=3-i$，则$\overline{z}=$
\begin{choices}
  \item $-1+2i$
  \item $1-2i$
  \item $3+2i$
  \item $3-2i$
\end{choices}
\begin{solution}
\textbf{答案：}$3+2i$

\textbf{分析：}由$z+i=3-i$得$z=3-2i$，所以$\overline{z}=3+2i$。
\end{solution}
\end{question}

\begin{question}[points=5]
\textbf{来源：}2016 年 全国I卷-理科，第 2 题\par
设$(1+i)x=1+yi$，其中$x$，$y$是实数，则$|x+yi|=$（ ）
\begin{choices}
  \item $1$
  \item $\sqrt{2}$
  \item $\sqrt{3}$
  \item $2$
\end{choices}
\begin{solution}
\textbf{答案：}B

\textbf{分析：}根据复数相等求出$x$，$y$的值，结合复数的模长公式进行计算即可．

\textbf{详解：}解：因为$(1+i)x=1+yi$，所以$x+xi=1+yi$，即$\begin{cases}x=1\\x=y\end{cases}$，解得$\begin{cases}x=1\\y=1\end{cases}$，即$|x+yi|=|1+i|=\sqrt{2}$，故选：B．
\end{solution}
\end{question}

\begin{question}[points=5]
\textbf{来源：}2016 年 全国I卷-文科，第 2 题\par
设$(1+2i)(a+i)$的实部与虚部相等，其中$a$为实数，则$a$等于
\begin{choices}
  \item $-3$
  \item $-2$
  \item $2$
  \item $3$
\end{choices}
\begin{solution}
\textbf{答案：}$-3$

\textbf{分析：}利用复数的乘法运算法则，通过复数相等的充要条件求解即可。

\textbf{详解：}$(1+2i)(a+i)=a-2+(2a+1)i$的实部与虚部相等，可得$a-2=2a+1$，解得$a=-3$。故选：$A$。
\end{solution}
\end{question}

\section*{2017 年}

\begin{question}[points=5]
\textbf{来源：}2017 年 全国III卷-理科，第 2 题\par
设复数$z$满足$(1+i)z=2i$，则$|z|=$（ ）
\begin{choices}
  \item $\frac{1}{2}$
  \item $\frac{\sqrt{2}}{2}$
  \item $\sqrt{2}$
  \item $2$
\end{choices}
\begin{solution}
\textbf{答案：}$C$

\textbf{分析：}利用复数的运算法则、模的计算公式即可得出。

\textbf{详解：}因为$(1+i)z=2i$，所以$(1-i)(1+i)z=2i(1-i)$，$z=i+1$。则$|z|=\sqrt{2}$。故选C。
\end{solution}
\end{question}

\begin{question}[points=5]
\textbf{来源：}2017 年 全国III卷-文科，第 2 题\par
复平面内表示复数$z=i(-2+i)$的点位于（            ）
\begin{choices}
  \item A．第一象限
  \item B．第二象限
  \item C．第三象限
  \item D．第四象限
\end{choices}
\begin{solution}
\textbf{答案：}C

\textbf{分析：}利用复数的运算法则、几何意义即可得出。

\textbf{详解：}$z=i(-2+i)=-2i-1$对应的点$(-1,-2)$位于第三象限。
\end{solution}
\end{question}

\begin{question}[points=5]
\textbf{来源：}2017 年 全国II卷-理科，第 1 题\par
$\frac{3+i}{1+i}=$（   ）
\begin{choices}
  \item $1+2i$
  \item $1-2i$
  \item $2+i$
  \item $2-i$
\end{choices}
\begin{solution}
\textbf{答案：}$D$

\textbf{分析：}分子和分母同时乘以分母的共轭复数，再利用虚数单位 $i$ 的幂运算性质，求出结果。

\textbf{详解：}$\frac{3+i}{1+i}=\frac{(3+i)(1-i)}{(1+i)(1-i)}=\frac{4-2i}{2}=2-i$，故选：$D$。
\end{solution}
\end{question}

\begin{question}[points=5]
\textbf{来源：}2017 年 全国II卷-文科，第 2 题\par
$(1+i)(2+i)=$
\begin{choices}
  \item $1-i$
  \item $1+3i$
  \item $3+i$
  \item $3+3i$
\end{choices}
\begin{solution}
\textbf{答案：}$1+3i$

\textbf{分析：}利用复数的运算法则即可得出.

\textbf{详解：}原式 $=2-1+3i=1+3i$. 故选B.
\end{solution}
\end{question}

\begin{question}[points=5]
\textbf{来源：}2017 年 全国I卷-理科，第 3 题\par
设有下面四个命题



$p_1$：若复数 $z$ 满足 $\frac{1}{z}\in\mathbb{R}$，则 $z\in\mathbb{R}$；



$p_2$：若复数 $z$ 满足 $z^2\in\mathbb{R}$，则 $z\in\mathbb{R}$；



$p_3$：若复数 $z_1$，$z_2$ 满足 $z_1z_2\in\mathbb{R}$，则 $z_1=\overline{z_2}$；



$p_4$：若复数 $z\in\mathbb{R}$，则 $\overline{z}\in\mathbb{R}$．

其中的真命题为（ ）
\begin{choices}
  \item $p_1$，$p_3$
  \item $p_1$，$p_4$
  \item $p_2$，$p_3$
  \item $p_2$，$p_4$
\end{choices}
\begin{solution}
\textbf{答案：}B

\textbf{分析：}根据复数的分类，有复数性质，逐一分析给定四个命题的真假，可得答案．

\textbf{详解：}若复数 $z$ 满足 $\frac{1}{z}\in\mathbb{R}$，则 $z\in\mathbb{R}$，故命题 $p_1$ 为真命题；$p_2$：复数 $z=i$ 满足 $z^2=-1\in\mathbb{R}$，则 $z\notin\mathbb{R}$，故命题 $p_2$ 为假命题；$p_3$：若复数 $z_1=i$，$z_2=2i$ 满足 $z_1z_2\in\mathbb{R}$，但 $z_1\neq \overline{z_2}$，故命题 $p_3$ 为假命题；$p_4$：若复数 $z\in\mathbb{R}$，则 $\overline{z}=z\in\mathbb{R}$，故命题 $p_4$ 为真命题．故选：B．
\end{solution}
\end{question}

\begin{question}[points=5]
\textbf{来源：}2017 年 全国I卷-文科，第 3 题\par
下列各式的运算结果为纯虚数的是（                ）
\begin{choices}
  \item $\text{i}(1+\text{i})^2$
  \item $\text{i}^2(1-\text{i})$
  \item $(1+\text{i})^2$
  \item $\text{i}(1+\text{i})$
\end{choices}
\begin{solution}
\textbf{答案：}C

\textbf{分析：}利用复数的运算法则、纯虚数的定义即可判断出结论。

\textbf{详解：}解：A．$\text{i}(1+\text{i})^2=\text{i}\cdot2\text{i}=-2$，是实数．B．$\text{i}^2(1-\text{i})=-1+\text{i}$，不是纯虚数．C．$(1+\text{i})^2=2\text{i}$为纯虚数．D．$\text{i}(1+\text{i})=\text{i}-1$不是纯虚数．故选：C。
\end{solution}
\end{question}

\section*{2018 年}

\begin{question}[points=5]
\textbf{来源：}2018 年 全国III卷-理科，第 2 题\par
$(1+i)(2-i)=$（ ）
\begin{choices}
  \item $-3-i$
  \item $-3+i$
  \item $3-i$
  \item $3+i$
\end{choices}
\begin{solution}
\textbf{答案：}D

\textbf{分析：}直接利用复数代数形式的乘除运算化简得答案．

\textbf{详解：}$(1+i)(2-i)=3+i$．
\end{solution}
\end{question}

\begin{question}[points=5]
\textbf{来源：}2018 年 全国III卷-文科，第 2 题\par
$(1+i)(2-i)=$
\begin{choices}
  \item $-3-i$
  \item $-3+i$
  \item $3-i$
  \item $3+i$
\end{choices}
\begin{solution}
\textbf{答案：}D

\textbf{分析：}直接利用复数代数形式的乘除运算化简得答案．

\textbf{详解：}$(1+i)(2-i)=3+i$．故选：D．
\end{solution}
\end{question}

\begin{question}[points=5]
\textbf{来源：}2018 年 全国II卷-理科，第 1 题\par
$\frac{1+2i}{1-2i}=$（    ）
\begin{choices}
  \item $-\frac{4}{5}-\frac{3}{5}i$
  \item $-\frac{4}{5}+\frac{3}{5}i$
  \item $-\frac{3}{5}-\frac{4}{5}i$
  \item $-\frac{3}{5}+\frac{4}{5}i$
\end{choices}
\begin{solution}
\textbf{答案：}D

\textbf{分析：}利用复数的除法的运算法则化简求解即可．

\textbf{详解：}$\frac{1+2i}{1-2i}=\frac{(1+2i)^2}{(1-2i)(1+2i)}=\frac{1+4i+4i^2}{1-4i^2}=\frac{-3+4i}{5}=-\frac{3}{5}+\frac{4}{5}i$．故选D．
\end{solution}
\end{question}

\begin{question}[points=5]
\textbf{来源：}2018 年 全国II卷-文科，第 1 题\par
$i(2+3i)=$（ ）
\begin{choices}
  \item $3-2i$
  \item $3+2i$
  \item $-3-2i$
  \item $-3+2i$
\end{choices}
\begin{solution}
\textbf{答案：}$-3+2i$

\textbf{分析：}利用复数的代数形式的乘除运算法则直接求解。

\textbf{详解：}$i(2+3i)=2i+3i^2=-3+2i$。故选：D。
\end{solution}
\end{question}

\begin{question}[points=5]
\textbf{来源：}2018 年 全国I卷-理科，第 1 题\par
设$z=\frac{1-i}{1+i}+2i$，则$|z|=$（    ）
\begin{choices}
  \item $0$
  \item $\frac{1}{2}$
  \item $1$
  \item $\sqrt{2}$
\end{choices}
\begin{solution}
\textbf{答案：}C

\textbf{分析：}利用复数的代数形式的混合运算化简后，然后求解复数的模。

\textbf{详解：}$z=\frac{1-i}{1+i}+2i=\frac{(1-i)^2}{(1+i)(1-i)}+2i=\frac{-2i}{2}+2i=-i+2i=i$，则$|z|=1$。故选C。
\end{solution}
\end{question}

\begin{question}[points=5]
\textbf{来源：}2018 年 全国I卷-文科，第 2 题\par
设 $z=\frac{1-i}{1+i}+2i$，则 $|z|=$
\begin{choices}
  \item $0$
  \item $\frac{1}{2}$
  \item $1$
  \item $\sqrt{2}$
\end{choices}
\begin{solution}
\textbf{答案：}$1$

\textbf{分析：}利用复数的代数形式的混合运算化简后，然后求解复数的模。

\textbf{详解：}$z=\frac{1-i}{1+i}+2i=\frac{(1-i)^2}{(1+i)(1-i)}+2i=\frac{-2i}{2}+2i=-i+2i=i$，则 $|z|=1$。故选：$C$。
\end{solution}
\end{question}

\section*{2019 年}

\begin{question}[points=5]
\textbf{来源：}2019 年 全国III卷-理科，第 2 题\par
若 $z(1+i)=2i$，则 $z=$
\begin{choices}
  \item $-1-i$
  \item $-1+i$
  \item $1-i$
  \item $1+i$
\end{choices}
\begin{solution}
\textbf{答案：}D

\textbf{分析：}根据复数运算法则求解即可。

\textbf{详解：}$z = \dfrac{2i}{1+i} = \dfrac{2i(1-i)}{(1+i)(1-i)} = 1+i$。故选 D。
\end{solution}
\end{question}

\begin{question}[points=5]
\textbf{来源：}2019 年 全国III卷-文科，第 2 题\par
若$z(1+i)=2i$，则$z=$
\begin{choices}
  \item $-1-i$
  \item $-1+i$
  \item $1-i$
  \item $1+i$
\end{choices}
\begin{solution}
\textbf{答案：}D

\textbf{分析：}根据复数运算法则求解即可。

\textbf{详解：}$z=\frac{2i}{1+i}=\frac{2i(1-i)}{(1+i)(1-i)}=1+i$。故选D。
\end{solution}
\end{question}

\begin{question}[points=5]
\textbf{来源：}2019 年 全国II卷-理科，第 2 题\par
设 $z=-3+2i$，则在复平面内 $\overline{z}$ 对应的点位于
\begin{choices}
  \item 第一象限
  \item 第二象限
  \item 第三象限
  \item 第四象限
\end{choices}
\begin{solution}
\textbf{答案：}$C$

\textbf{分析：}本题考查复数的共轭复数和复数在复平面内的对应点位置，渗透了直观想象和数学运算素养．采取定义法，利用数形结合思想解题．

\textbf{详解：}由 $z=-3+2i$，得 $\overline{z}=-3-2i$，则 $\overline{z}=-3-2i$，对应点 $(-3,-2)$ 位于第三象限．故选 $C$．
\end{solution}
\end{question}

\begin{question}[points=5]
\textbf{来源：}2019 年 全国II卷-文科，第 2 题\par
设 $z=\mathrm{i}(2+\mathrm{i})$，则 $\overline{z}=$
\begin{choices}
  \item $1+2\mathrm{i}$
  \item $-1+2\mathrm{i}$
  \item $1-2\mathrm{i}$
  \item $-1-2\mathrm{i}$
\end{choices}
\begin{solution}
\textbf{答案：}$D$

\textbf{分析：}本题根据复数的乘法运算法则先求得 $z$，然后根据共轭复数的概念，写出 $\overline{z}$．

\textbf{详解：}$z=\mathrm{i}(2+\mathrm{i})=2\mathrm{i}+\mathrm{i}^2=-1+2\mathrm{i}$，所以 $\overline{z}=-1-2\mathrm{i}$，选 $D$．
\end{solution}
\end{question}

\begin{question}[points=5]
\textbf{来源：}2019 年 全国I卷-理科，第 2 题\par
设复数 $z$ 满足 $|z - i| = 1$, $z$ 在复平面内对应的点为 $(x, y)$, 则
\begin{choices}
  \item $(x + 1)^2 + y^2 = 1$
  \item $(x - 1)^2 + y^2 = 1$
  \item $x^2 + (y - 1)^2 = 1$
  \item $x^2 + (y + 1)^2 = 1$
\end{choices}
\begin{solution}
\textbf{答案：}$C$

\textbf{分析：}本题考点为复数的运算，为基础题目，难度偏易．此题可采用几何法，根据点 $(x, y)$ 和点 $(0, 1)$ 之间的距离为 $1$，可选正确答案 $C$．

\textbf{详解：}$z = x + yi$, $z - i = x + (y - 1)i$, $|z - i| = \sqrt{x^2 + (y - 1)^2} = 1$, 则 $x^2 + (y - 1)^2 = 1$．故选 $C$．
\end{solution}
\end{question}

\begin{question}[points=5]
\textbf{来源：}2019 年 全国I卷-文科，第 1 题\par
设 $z = \frac{3-i}{1+2i}$，则 $|z| =$
\begin{choices}
  \item $2$
  \item $\sqrt{3}$
  \item $\sqrt{2}$
  \item $1$
\end{choices}
\begin{solution}
\textbf{答案：}C

\textbf{分析：}先由复数的除法运算（分母实数化），求得 $z$，再求 $|z|$。

\textbf{详解：}因为 $z = \frac{3-i}{1+2i}$，所以 $z = \frac{(3-i)(1-2i)}{(1+2i)(1-2i)} = \frac{1}{5} - \frac{7}{5}i$，所以 $|z| = \sqrt{\left(\frac{1}{5}\right)^2 + \left(-\frac{7}{5}\right)^2} = \sqrt{2}$，故选 C。
\end{solution}
\end{question}

\section*{2020 年}

\begin{question}[points=5]
\textbf{来源：}2020 年 全国III卷-理科，第 2 题\par
复数 $\dfrac{1}{1-3i}$ 的虚部是
\begin{choices}
  \item $-\dfrac{3}{10}$
  \item $-\dfrac{1}{10}$
  \item $\dfrac{1}{10}$
  \item $\dfrac{3}{10}$
\end{choices}
\begin{solution}
\textbf{答案：}$D$

\textbf{分析：}利用复数的除法运算求出 $z$ 即可。

\textbf{详解：}因为 $z=\dfrac{1}{1-3i}=\dfrac{1+3i}{(1-3i)(1+3i)}=\dfrac{1}{10}+\dfrac{3}{10}i$，所以复数 $z=\dfrac{1}{1-3i}$ 的虚部为 $\dfrac{3}{10}$。故选 $D$。
\end{solution}
\end{question}

\begin{question}[points=5]
\textbf{来源：}2020 年 全国III卷-文科，第 2 题\par
若 $z(1+i)=1-i$，则 $z=$
\begin{choices}
  \item $1-i$
  \item $1+i$
  \item $-i$
  \item $i$
\end{choices}
\begin{solution}
\textbf{答案：}D

\textbf{分析：}先利用除法运算求得 $z$，再利用共轭复数的概念得到 $z$ 即可.

\textbf{详解：}因为 $z=\frac{1-i}{1+i}=\frac{(1-i)^2}{(1+i)(1-i)}=\frac{-2i}{2}=-i$，所以 $\bar{z}=i$. 故选：D.
\end{solution}
\end{question}

\begin{question}[points=5]
\textbf{来源：}2020 年 全国II卷-理科，第 15 题\par
设复数$z_1,z_2$满足$|z_1|=|z_2|=2$，$z_1+z_2=\sqrt3+\mathrm{i}$，则$|z_1-z_2|=$.
\fillin[{$2\sqrt3$}]
\begin{solution}
\textbf{答案：}$2\sqrt3$

\textbf{分析：}令$z_1=2\cos\theta+2\sin\theta\cdot\mathrm{i}$，$z_2=2\cos\alpha+2\sin\alpha\cdot\mathrm{i}$，根据复数相等可求得$\cos\theta\cos\alpha+\sin\theta\sin\alpha=-\frac12$，代入复数模长的公式中即可得到结果.

\textbf{详解：}$\because|z_1|=|z_2|=2$，可设$z_1=2\cos\theta+2\sin\theta\cdot\mathrm{i}$，$z_2=2\cos\alpha+2\sin\alpha\cdot\mathrm{i}$，$\therefore z_1+z_2=2(\cos\theta+\cos\alpha)+2(\sin\theta+\sin\alpha)\mathrm{i}=\sqrt3+\mathrm{i}$，$\therefore\begin{cases}2(\cos\theta+\cos\alpha)=\sqrt3\\2(\sin\theta+\sin\alpha)=1\end{cases}$，两式平方作和得：$4(2+2(\cos\theta\cos\alpha+\sin\theta\sin\alpha))=4$，化简得$\cos\theta\cos\alpha+\sin\theta\sin\alpha=-\frac12$，$\therefore|z_1-z_2|^2=4(\cos\theta-\cos\alpha)^2+4(\sin\theta-\sin\alpha)^2=8-8(\cos\theta\cos\alpha+\sin\theta\sin\alpha)=8+4=12$，$\therefore|z_1-z_2|=2\sqrt3$.
\end{solution}
\end{question}

\begin{question}[points=5]
\textbf{来源：}2020 年 全国II卷-文科，第 2 题\par
$(1-i)^4 =$（ ）
\begin{choices}
  \item $-4$
  \item $4$
  \item $-4i$
  \item $4i$
\end{choices}
\begin{solution}
\textbf{答案：}A

\textbf{分析：}根据指数幂的运算性质，结合复数的乘方运算性质进行求解。

\textbf{详解：}$(1-i)^4 = [(1-i)^2]^2 = (1-2i+i^2)^2 = (-2i)^2 = -4$。
故选：A。
\end{solution}
\end{question}

\begin{question}[points=5]
\textbf{来源：}2020 年 全国I卷-理科，第 1 题\par
若$z=1+i$，则$|z^2-2z|=$
\begin{choices}
  \item $0$
  \item $1$
  \item $\sqrt{2}$
  \item $2$
\end{choices}
\begin{solution}
\textbf{答案：}D

\textbf{详解：}由题意可得$z^2=(1+i)^2=2i$，则$z^2-2z=2i-2(1+i)=-2$，故$|z^2-2z|=|-2|=2$。
\end{solution}
\end{question}

\begin{question}[points=5]
\textbf{来源：}2020 年 全国I卷-文科，第 2 题\par
若 $z = 1 + 2i + i^3$, 则 $|z| =$
\begin{choices}
  \item $0$
  \item $1$
  \item $\sqrt{2}$
  \item $2$
\end{choices}
\begin{solution}
\textbf{答案：}$\sqrt{2}$

\textbf{分析：}先根据 $i^2 = -1$ 将 $z$ 化简, 再根据向量的模的计算公式即可求出.

\textbf{详解：}因为 $z = 1 + 2i + i^3 = 1 + 2i - i = 1 + i$, 所以 $|z| = \sqrt{1^2 + 1^2} = \sqrt{2}$.
\end{solution}
\end{question}

\section*{2021 年}

\begin{question}[points=4]
\textbf{来源：}2021 年 北京卷，第 2 题\par
在复平面内, 复数 $z$ 满足 $(1-i)z=2$, 则 $z=$ （ ）
\begin{choices}
  \item $2+i$
  \item $2-i$
  \item $1-i$
  \item $1+i$
\end{choices}
\begin{solution}
\textbf{答案：}D

\textbf{分析：}由题意利用复数的运算法则整理计算即可求得最终结果.

\textbf{详解：}由题意可得 $z=\frac{2}{1-i}=\frac{2(1+i)}{(1-i)(1+i)}=\frac{2(1+i)}{2}=1+i$.
\end{solution}
\end{question}

\begin{question}[points=5]
\textbf{来源：}2021 年 全国甲卷-理科，第 3 题\par
已知 $(1-i)^2 z = 3 + 2i$, 则 $z =$ ( )
\begin{choices}
  \item $-1-\frac{3}{2}i$
  \item $-1+\frac{3}{2}i$
  \item $-\frac{3}{2}+i$
  \item $-\frac{3}{2}-i$
\end{choices}
\begin{solution}
\textbf{答案：}B

\textbf{分析：}由已知得 $z = \frac{3+2i}{-2i}$，根据复数除法运算法则，即可求解．

\textbf{详解：}$(1-i)^2 = -2i$, 则 $z = \frac{3+2i}{-2i} = \frac{(3+2i)\cdot i}{-2i\cdot i} = \frac{-2+3i}{2} = -1+\frac{3}{2}i$．
\end{solution}
\end{question}

\begin{question}[points=5]
\textbf{来源：}2021 年 全国甲卷-文科，第 3 题\par
已知 $(1-i)^2 z = 3+2i$, 则 $z =$ ( )
\begin{choices}
  \item $-1-\frac{3}{2}i$
  \item $-1+\frac{3}{2}i$
  \item $-\frac{3}{2}+i$
  \item $-\frac{3}{2}-i$
\end{choices}
\begin{solution}
\textbf{答案：}B

\textbf{分析：}由已知得 $z = \frac{3+2i}{-2i}$, 根据复数除法运算法则，即可求解.

\textbf{详解：}$(1-i)^2 z = -2i z = 3+2i$, $z = \frac{3+2i}{-2i} = \frac{(3+2i)\cdot i}{-2i \cdot i} = \frac{-2+3i}{2} = -1 + \frac{3}{2}i$.
\end{solution}
\end{question}

\begin{question}[points=5]
\textbf{来源：}2021 年 全国乙卷-理科，第 1 题\par
设$2(z+\overline{z})+3(z-\overline{z})=4+6i$，则$z=$（  ）
\begin{choices}
  \item $1-2i$
  \item $1+2i$
  \item $1+i$
  \item $1-i$
\end{choices}
\begin{solution}
\textbf{答案：}$1+i$

\textbf{详解：}设$z=a+bi$，则$\overline{z}=a-bi$，$2(z+\overline{z})+3(z-\overline{z})=4a+6bi=4+6i$，所以$a=1$，$b=1$，所以$z=1+i$。故选C。
\end{solution}
\end{question}

\begin{question}[points=5]
\textbf{来源：}2021 年 全国乙卷-文科，第 2 题\par
设 $iz = 4 + 3i$ ，则 $z =$ （ ）
\begin{choices}
  \item $-3-4i$
  \item $-3+4i$
  \item $3-4i$
  \item $3+4i$
\end{choices}
\begin{solution}
\textbf{答案：}C
\end{solution}
\end{question}

\begin{question}[points=4]
\textbf{来源：}2021 年 上海春季卷，第 2 题\par
已知 $z=1-3i$，则 $|z-i|=$  .
\fillin[{$\sqrt{5}$}]
\begin{solution}
\textbf{答案：}$\sqrt{5}$

\textbf{分析：}由已知求得 $z-i$，再由复数模的计算公式求解．

\textbf{详解：}因为 $z=1-3i$，所以 $z-i=1+3i-i=1+2i$，则 $|z-i|=|1+2i|=\sqrt{1^2+2^2}=\sqrt{5}$．
\end{solution}
\end{question}

\begin{question}[points=4]
\textbf{来源：}2021 年 上海卷，第 1 题\par
已知$z_1=1+i$, $z_2=2+3i$（其中$i$为虚数单位），则$z_1+z_2=$．
\fillin[{$3+4i$}]
\begin{solution}
\textbf{答案：}$3+4i$

\textbf{分析：}复数实部和虚部分别相加

\textbf{详解：}$z_1+z_2=3+4i$
\end{solution}
\end{question}

\begin{question}[points=5]
\textbf{来源：}2021 年 天津卷，第 10 题\par
$i$ 是虚数单位，复数 $\frac{9+2i}{2+i} =$ .
\fillin[{$4-i$}]
\begin{solution}
\textbf{答案：}$4-i$

\textbf{分析：}利用复数的除法化简可得结果.

\textbf{详解：}$\frac{9+2i}{2+i} = \frac{(9+2i)(2-i)}{(2+i)(2-i)} = \frac{20-5i}{5} = 4-i$.
\end{solution}
\end{question}

\begin{question}[points=5]
\textbf{来源：}2021 年 新高考II卷，第 1 题\par
复数$\frac{2-i}{1-3i}$在复平面内对应的点所在的象限为（ ）
\begin{choices}
  \item A. 第一象限
  \item B. 第二象限
  \item C. 第三象限
  \item D. 第四象限
\end{choices}
\begin{solution}
\textbf{答案：}A

\textbf{分析：}利用复数的除法可化简$\frac{2-i}{1-3i}$，从而可求对应的点的位置.

\textbf{详解：}$\frac{2-i}{1-3i}=\frac{(2-i)(1+3i)}{10}=\frac{5+5i}{10}=\frac{1+i}{2}$，所以该复数对应的点为$(\frac{1}{2},\frac{1}{2})$，该点在第一象限，故选：A.
\end{solution}
\end{question}

\begin{question}[points=5]
\textbf{来源：}2021 年 新高考I卷，第 2 题\par
已知 $z = 2 - i$, 则 $z(\overline{z} + i) =$ ( )
\begin{choices}
  \item $6 - 2i$
  \item $4 - 2i$
  \item $6 + 2i$
  \item $4 + 2i$
\end{choices}
\begin{solution}
\textbf{答案：}$6 + 2i$

\textbf{分析：}利用复数的乘法和共轭复数的定义可求得结果.

\textbf{详解：}因为 $z = 2 - i$, 故 $\overline{z} = 2 + i$, 故 $z(\overline{z} + i) = (2 - i)(2 + 2i) = 6 + 2i$, 故选C.
\end{solution}
\end{question}

\begin{question}[points=4]
\textbf{来源：}2021 年 浙江卷，第 2 题\par
已知 $a \in \mathbb{R}$, $(1+ai)i=3+i$, (i为虚数单位), 则 $a=$ （   ）
\begin{choices}
  \item $-1$
  \item $1$
  \item $-3$
  \item $3$
\end{choices}
\begin{solution}
\textbf{答案：}C

\textbf{分析：}首先计算左侧的结果，然后结合复数相等的充分必要条件即可求得实数 $a$ 的值.

\textbf{详解：}$(1+ai)i = i + ai^2 = i - a = -a + i = 3 + i$, 所以 $-a=3$, 故 $a=-3$.
\end{solution}
\end{question}

\section*{2022 年}

\begin{question}[points=4]
\textbf{来源：}2022 年 北京卷，第 2 题\par
若复数 $z$ 满足 $i \cdot z = 3 - 4i$，则 $|z| =$
\begin{choices}
  \item $1$
  \item $5$
  \item $7$
  \item $25$
\end{choices}
\begin{solution}
\textbf{答案：}B

\textbf{分析：}利用复数四则运算，先求出 $z$，再计算复数的模。

\textbf{详解：}由题意有 $z = \dfrac{3-4i}{i} = \dfrac{(3-4i)(-i)}{i \cdot (-i)} = -4-3i$，故 $|z| = \sqrt{(-4)^2 + (-3)^2} = 5$。故选B。
\end{solution}
\end{question}

\begin{question}[points=5]
\textbf{来源：}2022 年 全国甲卷-理科，第 1 题\par
若 $z = -1 + \sqrt{3}i$，则 $\frac{z}{z\bar{z} - 1} = $（ ）
\begin{choices}
  \item $-1+\sqrt{3}i$
  \item $-1-\sqrt{3}i$
  \item $-\frac{1}{3}+\frac{\sqrt{3}}{3}i$
  \item $-\frac{1}{3}-\frac{\sqrt{3}}{3}i$
\end{choices}
\begin{solution}
\textbf{答案：}C

\textbf{分析：}由共轭复数的概念及复数的运算即可得解.

\textbf{详解：}$\bar{z} = -1 - \sqrt{3}i$, $z\bar{z} = (-1+\sqrt{3}i)(-1-\sqrt{3}i) = 1+3 = 4$. 则 $\frac{z}{z\bar{z}-1} = \frac{-1+\sqrt{3}i}{3} = -\frac{1}{3} + \frac{\sqrt{3}}{3}i$. 故选C.
\end{solution}
\end{question}

\begin{question}[points=5]
\textbf{来源：}2022 年 全国甲卷-文科，第 3 题\par
若 $z=1+\mathrm{i}$，则 $|\mathrm{i}z+3\overline{z}|=$
\begin{choices}
  \item $4\sqrt{5}$
  \item $4\sqrt{2}$
  \item $2\sqrt{5}$
  \item $2\sqrt{2}$
\end{choices}
\begin{solution}
\textbf{答案：}D
\end{solution}
\end{question}

\begin{question}[points=5]
\textbf{来源：}2022 年 全国乙卷-理科，第 2 题\par
已知$z = 1 - 2\mathrm{i}$，且$z + a\bar{z} + b = 0$，其中$a$，$b$为实数，则（\qquad）
\begin{choices}
  \item $a = 1, b = -2$
  \item $a = -1, b = 2$
  \item $a = 1, b = 2$
  \item $a = -1, b = -2$
\end{choices}
\begin{solution}
\textbf{答案：}A

\textbf{分析：}先算出$\bar{z}$,再代入计算,实部与虚部都为零解方程组即可

\textbf{详解：}$\bar{z} = 1 + 2\mathrm{i}$\\$z + a\bar{z} + b = 1 - 2\mathrm{i} + a(1 + 2\mathrm{i}) + b = (1 + a + b) + (2a - 2)\mathrm{i}$由$z + a\bar{z} + b = 0$,得$\begin{cases} 1 + a + b = 0 \\ 2a - 2 = 0 \end{cases}$,即$\begin{cases} a = 1 \\ b = -2 \end{cases}$
\end{solution}
\end{question}

\begin{question}[points=5]
\textbf{来源：}2022 年 全国乙卷-文科，第 2 题\par
设 $(1+2i)a + b = 2i$, 其中 $a, b$ 为实数，则
\begin{choices}
  \item $a=1, b=-1$
  \item $a=1, b=1$
  \item $a=-1, b=1$
  \item $a=-1, b=-1$
\end{choices}
\begin{solution}
\textbf{答案：}A

\textbf{分析：}根据复数代数形式的运算法则以及复数相等的概念即可解出.

\textbf{详解：}因为 $a,b\in\mathbb{R}$, $(a+b)+2ai = 2i$, 所以 $a+b=0, 2a=2$, 解得 $a=1, b=-1$.
\end{solution}
\end{question}

\begin{question}[points=4]
\textbf{来源：}2022 年 上海卷，第 3 题\par
\begin{center}

\begin{tabular}{ccccc}

已知$a \in \mathbb{R}$，行列式$\begin{vmatrix} a & 1 \\ 3 & 2 \end{vmatrix}$的值与行列式$\begin{vmatrix} a & 0 \\ 4 & 1 \end{vmatrix}$的值相等，则$a=$ ．

\end{tabular}

\end{center}
\fillin[{$3$}]
\begin{solution}
\textbf{答案：}$3$

\textbf{分析：}根据行列式所表示的值求解即可．

\textbf{详解：}解：因为$\begin{vmatrix} a & 1 \\ 3 & 2 \end{vmatrix}=2a-3$，$\begin{vmatrix} a & 0 \\ 4 & 1 \end{vmatrix}=a$，所以$2a-3=a$，解得$a=3$．故答案为：$3$．
\end{solution}
\end{question}

\begin{question}[points=5]
\textbf{来源：}2022 年 天津卷，第 10 题\par
已知 $i$ 是虚数单位，化简 $\frac{11-3i}{1+2i}$ 的结果为．
\fillin[{$1-5i$}]
\begin{solution}
\textbf{答案：}$1-5i$

\textbf{分析：}$\frac{11-3i}{1+2i}=\frac{(11-3i)(1-2i)}{(1+2i)(1-2i)}=\frac{11-6-25i}{5}=1-5i$.
\end{solution}
\end{question}

\begin{question}[points=5]
\textbf{来源：}2022 年 新高考II卷，第 2 题\par
$(2+2i)(1-2i)=$（ ）
\begin{choices}
  \item $-2+4i$
  \item $-2-4i$
  \item $6+2i$
  \item $6-2i$
\end{choices}
\begin{solution}
\textbf{答案：}$6-2i$

\textbf{分析：}利用复数的乘法可求 $(2+2i)(1-2i)$.

\textbf{详解：}$(2+2i)(1-2i)=2+4-4i+2i=6-2i$，故选D.
\end{solution}
\end{question}

\begin{question}[points=5]
\textbf{来源：}2022 年 新高考I卷，第 2 题\par
若 $\mathrm{i}(1 - z) = 1$, 则 $z + \bar{z} =$
\begin{choices}
  \item $-2$
  \item $-1$
  \item $1$
  \item $2$
\end{choices}
\begin{solution}
\textbf{答案：}D

\textbf{分析：}利用复数的除法可求 $z$, 从而可求 $z + \bar{z}$.

\textbf{详解：}由题设有 $1 - z = \frac{1}{\mathrm{i}} = \frac{\mathrm{i}}{\mathrm{i}^2} = -\mathrm{i}$, 故 $z = 1 + \mathrm{i}$, 故 $z + \bar{z} = (1 + \mathrm{i}) + (1 - \mathrm{i}) = 2$.
\end{solution}
\end{question}

\begin{question}[points=4]
\textbf{来源：}2022 年 浙江卷，第 2 题\par
已知 $a,b \in \mathbb{R}, a+3i = (b+i)i$ （$i$ 为虚数单位），则（ ）
\begin{choices}
  \item $a = 1, b = -3$
  \item $a = -1, b = 3$
  \item $a = -1, b = -3$
  \item $a = 1, b = 3$
\end{choices}
\begin{solution}
\textbf{答案：}B

\textbf{分析：}利用复数相等的条件可求 $a,b$.

\textbf{详解：}由 $a+3i = (b+i)i = -1 + bi$，得 $a = -1, b = 3$，故选B.
\end{solution}
\end{question}

\section*{2023 年}

\begin{question}[points=4]
\textbf{来源：}2023 年 北京卷，第 2 题\par
在复平面内，复数 $z$ 对应的点的坐标是 $(-1, \sqrt{3})$ ，则 $z$ 的共轭复数 $\bar{z} =$ （ ）
\begin{choices}
  \item $1 + \sqrt{3}i$
  \item $1 - \sqrt{3}i$
  \item $-1 + \sqrt{3}i$
  \item $-1 - \sqrt{3}i$
\end{choices}
\begin{solution}
\textbf{答案：}D

\textbf{分析：}根据复数的几何意义先求出复数 $z$ ，然后利用共轭复数的定义计算.

\textbf{详解：}$z$ 在复平面对应的点是 $(-1, \sqrt{3})$ ，根据复数的几何意义， $z = -1 + \sqrt{3}i$ ，由共轭复数的定义可知， $\bar{z} = -1 - \sqrt{3}i$ .
\end{solution}
\end{question}

\begin{question}[points=5]
\textbf{来源：}2023 年 全国甲卷-理科，第 2 题\par
设 $a \in R$，$(a + i)(1 - ai) = 2$，则 $a = $ （       ）
\begin{choices}
  \item $-1$
  \item $0$
  \item $1$
  \item $2$
\end{choices}
\begin{solution}
\textbf{答案：}C

\textbf{分析：}根据复数的代数运算以及复数相等即可解出．

\textbf{详解：}因为 $(a + i)(1 - ai) = a - a^2 i + i + a = 2a + (1 - a^2)i = 2$，所以 $\begin{cases} 2a = 2 \\ 1 - a^2 = 0 \end{cases}$，解得：$a = 1$．
\end{solution}
\end{question}

\begin{question}[points=5]
\textbf{来源：}2023 年 全国甲卷-文科，第 2 题\par
$\frac{5(1+i^3)}{(2+i)(2-i)} =$（    ）
\begin{choices}
  \item $-1$
  \item $1$
  \item $1-i$
  \item $1+i$
\end{choices}
\begin{solution}
\textbf{答案：}C

\textbf{分析：}利用复数的四则运算求解即可.

\textbf{详解：}$\frac{5(1+i^3)}{(2+i)(2-i)} = \frac{5(1-i)}{5} = 1-i$，故选C.
\end{solution}
\end{question}

\begin{question}[points=5]
\textbf{来源：}2023 年 全国乙卷-理科，第 1 题\par
设 $z=\frac{2+i}{1+i^{2}+i^{5}}$，则 $\overline{z}=$（ ）
\begin{choices}
  \item $1-2i$
  \item $1+2i$
  \item $2-i$
  \item $2+i$
\end{choices}
\begin{solution}
\textbf{答案：}$1+2i$

\textbf{分析：}由题意首先计算复数 $z$ 的值，然后利用共轭复数的定义确定其共轭复数即可.

\textbf{详解：}由题意可得 $z=\frac{2+i}{1+i^{2}+i^{5}}=\frac{2+i}{1-1+i}=\frac{2+i}{i}=\frac{i(2+i)}{i^{2}}=\frac{2i-1}{-1}=1-2i$，则 $\overline{z}=1+2i$.
\end{solution}
\end{question}

\begin{question}[points=5]
\textbf{来源：}2023 年 全国乙卷-文科，第 1 题\par
$2 + i^2 + 2i^3 =$（   ）
\begin{choices}
  \item $1$
  \item $2$
  \item $\sqrt{5}$
  \item $5$
\end{choices}
\begin{solution}
\textbf{答案：}$\sqrt{5}$

\textbf{分析：}由题意首先化简$2+i^2+2i^3$，然后计算其模即可.

\textbf{详解：}由题意可得$2+i^2+2i^3 = 2-1-2i = 1-2i$，则$|2+i^2+2i^3| = |1-2i| = \sqrt{1^2+(-2)^2} = \sqrt{5}$. 故选C.
\end{solution}
\end{question}

\begin{question}[points=4]
\textbf{来源：}2023 年 上海卷，第 2 题\par
已知 $a+2i=b+i$（$a,b\in \mathbb{R}$），求 $a+b=$
\fillin[{$4$}]
\begin{solution}
\textbf{答案：}$4$
\end{solution}
\end{question}

\begin{question}[points=5]
\textbf{来源：}2023 年 天津卷，第 10 题\par
已知 $i$ 是虚数单位，化简 $\dfrac{5 + 14i}{2 + 3i}$ 的结果为．
\fillin[{$4 + i$}]
\begin{solution}
\textbf{答案：}$4 + i$

\textbf{分析：}分子分母同时乘以 $2 - 3i$.

\textbf{详解：}$\dfrac{5 + 14i}{2 + 3i} = \dfrac{(5+14i)(2-3i)}{(2+3i)(2-3i)} = \dfrac{52 + 13i}{13} = 4 + i$.
\end{solution}
\end{question}

\begin{question}[points=5]
\textbf{来源：}2023 年 新高考II卷，第 1 题\par
在复平面内，$(1+3i)(3-i)$对应的点位于（ \quad ）
\begin{choices}
  \item A. 第一象限
  \item B. 第二象限
  \item C. 第三象限
  \item D. 第四象限
\end{choices}
\begin{solution}
\textbf{答案：}A

\textbf{分析：}根据复数的乘法结合复数的几何意义分析判断.

\textbf{详解：}因为$(1+3i)(3-i)=3+8i-3i^2=6+8i$，则所求复数对应的点为$(6,8)$，位于第一象限.
\end{solution}
\end{question}

\begin{question}[points=5]
\textbf{来源：}2023 年 新高考I卷，第 2 题\par
已知 $z=\frac{1-i}{2+2i}$，则 $z-\bar{z}=$（ ）
\begin{choices}
  \item $-i$
  \item $i$
  \item $0$
  \item $1$
\end{choices}
\begin{solution}
\textbf{答案：}$-i$

\textbf{分析：}根据复数的除法运算求出 $z$，再由共轭复数的概念得到 $\bar{z}$，从而解出。

\textbf{详解：}因为 $z=\frac{1-i}{2+2i}=\frac{(1-i)(1-i)}{2(1+i)(1-i)}=\frac{-2i}{4}=-\frac{1}{2}i$，所以 $\bar{z}=\frac{1}{2}i$，即 $z-\bar{z}=-i$。
\end{solution}
\end{question}

\section*{2024 年}

\begin{question}[points=4]
\textbf{来源：}2024 年 北京卷，第 2 题\par
已知 $\dfrac{z}{i} = i-1$，则 $z=$ （ ）
\begin{choices}
  \item $1-i$
  \item $-i$
  \item $-1-i$
  \item $1$
\end{choices}
\begin{solution}
\textbf{答案：}C

\textbf{分析：}直接根据复数乘法即可得到答案。

\textbf{详解：}由题意得 $z = i(i-1) = -1-i$，故选：C。
\end{solution}
\end{question}

\begin{question}[points=5]
\textbf{来源：}2024 年 全国甲卷-理科，第 1 题\par
设 $z=5+\mathrm{i}$，则 $\mathrm{i}(z+\overline{z})=$ （   ）
\begin{choices}
  \item $10\mathrm{i}$
  \item $2\mathrm{i}$
  \item $10$
  \item $-2$
\end{choices}
\begin{solution}
\textbf{答案：}A

\textbf{分析：}结合共轭复数与复数的基本运算直接求解。

\textbf{详解：}由 $z=5+\mathrm{i}\Rightarrow\overline{z}=5-\mathrm{i},z+\overline{z}=10$，则 $\mathrm{i}(z+\overline{z})=10\mathrm{i}$。
\end{solution}
\end{question}

\begin{question}[points=5]
\textbf{来源：}2024 年 全国甲卷-文科，第 2 题\par
设 $z = \sqrt{2}i$，则 $z \cdot \overline{z} =$（ ）
\begin{choices}
  \item $-i$
  \item $1$
  \item $-1$
  \item $2$
\end{choices}
\begin{solution}
\textbf{答案：}D

\textbf{分析：}先根据共轭复数的定义写出 $\overline{z}$，然后根据复数的乘法计算.

\textbf{详解：}依题意得，$\overline{z} = -\sqrt{2}i$，故 $z \cdot \overline{z} = (-\sqrt{2}i)(\sqrt{2}i) = 2$.
\end{solution}
\end{question}

\begin{question}[points=5]
\textbf{来源：}2024 年 上海卷，第 9 题\par
已知虚数$z$，其实部为1，且$z+\frac{2}{z}=m(m\in\mathbb{R})$，则实数$m$为.
\fillin[{$2$}]
\begin{solution}
\textbf{答案：}$2$

\textbf{分析：}设$z=1+bi$，直接根据复数的除法运算，再根据复数分类即可得到答案.

\textbf{详解：}设$z=1+bi$，$b\in\mathbb{R}$且$b\neq0$．则$z+\frac{2}{z}=1+bi+\frac{2}{1+bi}=\left(\frac{b^2+3}{1+b^2}\right)+\left(\frac{b^3-b}{1+b^2}\right)i=m$，$\because m\in\mathbb{R}$，$\therefore \begin{cases} \frac{b^2+3}{1+b^2}=m \\ \frac{b^3-b}{1+b^2}=0 \end{cases}$，解得$m=2$，故答案为：2.
\end{solution}
\end{question}

\begin{question}[points=5]
\textbf{来源：}2024 年 天津卷，第 10 题\par
已知 $i$ 是虚数单位，复数 $(\sqrt{5}+i)(\sqrt{5}-2i)=$ ．
\fillin[{$7-\sqrt{5}i$}]
\begin{solution}
\textbf{答案：}$7-\sqrt{5}i$

\textbf{分析：}直接利用复数乘法法则计算。

\textbf{详解：}$(\sqrt{5}+i)(\sqrt{5}-2i)=5-2\sqrt{5}i+\sqrt{5}i-2i^2=7-\sqrt{5}i$。
\end{solution}
\end{question}

\begin{question}[points=5]
\textbf{来源：}2024 年 新高考II卷，第 1 题\par
已知$z = -1 - i$，则$|z| =$（ ）
\begin{choices}
  \item $0$
  \item $1$
  \item $\sqrt{2}$
  \item $2$
\end{choices}
\begin{solution}
\textbf{答案：}$C$

\textbf{分析：}由复数模的计算公式直接计算即可.

\textbf{详解：}若$z = -1 - i$，则$|z| = \sqrt{(-1)^2 + (-1)^2} = \sqrt{2}$.
\end{solution}
\end{question}

\begin{question}[points=5]
\textbf{来源：}2024 年 新高考I卷，第 2 题\par
若 $\dfrac{z}{z-1}=1+\mathrm{i}$，则 $z=$（ ）
\begin{choices}
  \item $-1-\mathrm{i}$
  \item $-1+\mathrm{i}$
  \item $1-\mathrm{i}$
  \item $1+\mathrm{i}$
\end{choices}
\begin{solution}
\textbf{答案：}C

\textbf{分析：}由复数四则运算法则直接运算即可求解.

\textbf{详解：}因为 $\dfrac{z}{z-1}=1+\dfrac{1}{z-1}=1+\mathrm{i}$，所以 $\dfrac{1}{z-1}=\mathrm{i}$，故 $z=1+\dfrac{1}{\mathrm{i}}=1-\mathrm{i}$.
\end{solution}
\end{question}

\section*{2025 年}

\begin{question}[points=5]
\textbf{来源：}2025 年 全国II卷，第 2 题\par
已知 $z=1+\mathrm{i}$，则 $\frac{1}{z-1}=$（ ）
\begin{choices}
  \item $-\mathrm{i}$
  \item $\mathrm{i}$
  \item $-1$
  \item $1$
\end{choices}
\begin{solution}
\textbf{答案：}$A$

\textbf{分析：}由复数除法即可求解.

\textbf{详解：}因为 $z=1+\mathrm{i}$，所以 $\frac{1}{z-1}=\frac{1}{1+\mathrm{i}-1}=\frac{1}{\mathrm{i}}=-\mathrm{i}$.
\end{solution}
\end{question}

\begin{question}[points=5]
\textbf{来源：}2025 年 全国I卷，第 1 题\par
$(1+5i)i$ 的虚部为（ ）
\begin{choices}
  \item $-1$
  \item $0$
  \item $1$
  \item $6$
\end{choices}
\begin{solution}
\textbf{答案：}C

\textbf{分析：}根据复数代数形式的运算法则以及虚部的定义即可求出.

\textbf{详解：}因为 $(1+5i)i = i + 5i^2 = -5 + i$，所以其虚部为 $1$，故选：C.
\end{solution}
\end{question}

\begin{question}[points=5]
\textbf{来源：}2025 年 上海卷，第 10 题\par
已知复数 $z$ 满足 $z^2 = (\overline{z})^2, |z| \leq 1$，则 $|z - 2 - 3i|$ 的最小值是．
\fillin[{$2\sqrt{2}$}]
\begin{solution}
\textbf{答案：}$2\sqrt{2}$

\textbf{分析：}先设 $z = a + bi$，利用复数的乘方运算及概念确定 $ab = 0$，再根据复数的几何意义数形结合计算即可.

\textbf{详解：}设 $z = a + bi\,(a,b \in \mathbf{R})$，则 $\overline{z} = a - bi$，由题意可知 $z^2 = a^2 + 2abi - b^2 = \overline{z}^2 = a^2 - 2abi - b^2$，则 $ab = 0$，又 $|z| = \sqrt{a^2 + b^2} \leq 1$，由复数的几何意义知 $z$ 在复平面内对应的点 $Z(a,b)$ 在单位圆内部（含边界）的坐标轴上运动，如图所示即线段 $AB, CD$ 上运动，设 $E(2,3)$，则 $|z-2-3i| = ZE$，由图象可知 $BE = \sqrt{10} > CE = 2\sqrt{2}$，所以 $ZE_{\min} = 2\sqrt{2}$ .
\end{solution}
\end{question}

\begin{question}[points=5]
\textbf{来源：}2025 年 天津卷，第 10 题\par
已知$i$是虚数单位，则$\left| \dfrac{3+i}{i} \right| = $.
\fillin[{$\sqrt{10}$}]
\begin{solution}
\textbf{答案：}$\sqrt{10}$

\textbf{分析：}先由复数除法运算化简$\dfrac{3+i}{i}$, 再由复数模长公式即可计算求解.

\textbf{详解：}先由题得$\dfrac{3+i}{i} = -i(3+i) = 1 - 3i$, 所以$\left| \dfrac{3+i}{i} \right| = \sqrt{1^2 + (-3)^2} = \sqrt{10}$. 故答案为$\sqrt{10}$.
\end{solution}
\end{question}

